\documentclass[aspectratio=169, 10pt]{beamer}
\usetheme{metropolis}

% --- Edinburgh colours ----------------------------------------
\definecolor{UoERed}{HTML}{7A2318}
\definecolor{UoEGold}{HTML}{B8860B}
\definecolor{CodeBG}{HTML}{F7F3ED}
\definecolor{CodeFrame}{HTML}{D5C9B1}

\setbeamercolor{palette primary}{bg=UoERed, fg=white}
\setbeamercolor{frametitle}{bg=UoERed, fg=white}
\setbeamercolor{title separator}{fg=UoEGold}
\setbeamercolor{progress bar}{fg=UoEGold, bg=UoEGold!20}
\setbeamercolor{alerted text}{fg=UoERed}
\setbeamercolor{block title}{bg=UoERed!15, fg=UoERed}
\setbeamercolor{block body}{bg=UoERed!5}

% --- Packages -------------------------------------------------
\usepackage[T1]{fontenc}
\usepackage{lmodern}
\usepackage{amsmath, amssymb, mathtools}
\usepackage{listings}
\usepackage{graphicx, booktabs}
\usepackage{tikz}
\usetikzlibrary{arrows.meta, positioning, calc, decorations.pathreplacing}
\usepackage{hyperref}
\hypersetup{colorlinks=true, linkcolor=UoERed, urlcolor=UoERed!70!black}
\setbeamertemplate{navigation symbols}{}

\lstset{
  language=Python,
  basicstyle=\ttfamily\scriptsize,
  keywordstyle=\color{UoERed}\bfseries,
  commentstyle=\color{gray}\itshape,
  stringstyle=\color{UoEGold},
  backgroundcolor=\color{CodeBG},
  frame=single,
  rulecolor=\color{CodeFrame},
  numbers=none, breaklines=true, showstringspaces=false, tabsize=4,
  xleftmargin=4pt, xrightmargin=4pt, aboveskip=6pt, belowskip=4pt,
}

\AtBeginSection[]{
  \begin{frame}
    \vfill\centering
    \usebeamerfont{title}\insertsectionhead\par
    \vfill
  \end{frame}
}

\title{Week 7 --- Linear Programming}
\subtitle{The Simplex Method \& Economic Models}
\author{Dr Juan Zurita}
\institute{Optimisation \& Mathematical Methods for Economics\\
The University of Edinburgh~$\cdot$~School of Economics}
\date{}

\begin{document}
\maketitle

\begin{frame}{Linear Programs}
$\min\;\mathbf{c}^T\mathbf{x}$ s.t.\ $A\mathbf{x} \leq \mathbf{b}$, $\mathbf{x} \geq 0$\\\bigskip\begin{block}{Fundamental Theorem of LP}If an LP has an optimal solution, there is one at a \textbf{vertex} of the feasible polyhedron.\end{block}
\end{frame}

\begin{frame}{The Simplex Method}
Dantzig (1947): move from vertex to vertex, improving the objective at each step.\\\bigskip\begin{itemize}\item Worst case: exponential; in practice: very fast\item The most successful algorithm in operations research\end{itemize}
\end{frame}

\begin{frame}{LP Duality \& Shadow Prices}
Every LP has a \textbf{dual}. Optimal primal value = optimal dual value.\\\bigskip Dual variables = \textbf{shadow prices}: how much the objective improves per unit relaxation of each constraint.\\\bigskip In economics: the marginal value of an additional unit of a scarce resource.
\end{frame}

\begin{frame}{Applications}
\begin{itemize}\item \textbf{Input--output models:} Leontief's production planning\item \textbf{Transportation:} minimise shipping cost across a network\item \textbf{Diet problem:} minimise food cost meeting nutrition targets\item \textbf{Production planning:} allocate resources across products\end{itemize}\par\bigskip In Python: \texttt{scipy.optimize.linprog} or \texttt{PuLP}.
\end{frame}

\begin{frame}{Summary}
\begin{itemize}\item LP: optimise linear objective over linear constraints\item Simplex moves along vertices; solution is always at a corner\item Duality provides shadow prices for economic interpretation\item \textbf{Next week:} dynamic optimisation --- optimising over time\end{itemize}
\end{frame}
\end{document}
